\begin{pastexam}[label={px:ncku-115-prog-17}]{成大資訊聯招 115 學年度 演算法問答第 4 題}
（問答）

[10\%] You are optimizing the core engine of an image editing tool called ``Smart Lasso''. The tool helps
users extract objects by automatically finding the best boundary between a \textbf{Start Pixel (S)} and a
\textbf{Target Pixel (T)}.

The image is represented as a grid of pixels. Each pixel has an associated ``Energy Cost'' (representing
high-frequency noise or weak edges). Here are some rules:
\begin{itemize}[label=$\blacksquare$]
\item \textbf{Movement Rule}: From any pixel, you can move to adjacent pixels (Up, Down, Left, Right).
Diagonal movement is not allowed.
\item \textbf{Path Cost}: The sum of the Energy Costs of all pixels visited on the path (excluding S, but
including T).
\end{itemize}
Energy cost matrix is shown as below (assume the cost of entering T is 0):
\[
\begin{bmatrix}
S & 4 & 5 & 9 & 8\\
7 & 9 & 8 & 6 & 9\\
3 & 2 & 4 & 1 & 5\\
6 & 5 & 6 & 7 & T
\end{bmatrix}
\]
In the following, we will use (row, col) as coordinate representation. Top-left is (0, 0).

If two neighbors have the same cost, prioritize the direction in this order:\\
Right $\to$ Down $\to$ Up $\to$ Left

\underline{\textbf{Questions}}:
\begin{enumerate}[label=(\arabic*)]
\item {[4\%]} A junior developer implemented a ``High-Speed Mode'' for the lasso. The logic is simple:
\begin{itemize}[label=$\triangleright$]
\item ``At every step, the cursor simply snaps to the immediate neighbor with the lowest energy cost to
get closer to the target. It does not plan ahead.''
\end{itemize}
Please \textbf{trace the path} generated by this method from \textbf{S} to \textbf{T} and
\textbf{calculate the total cost.}
\begin{itemize}
\item Path Sequence: \underline{\hspace{4.5cm}} (e.g., (0,0) $\to$ (0,1) $\to$ \dots)
\item Total Path Cost: \underline{\hspace{4.5cm}}
\end{itemize}
\item {[4\%]} Users have complained that the ``High-Speed Mode'' (from Q 4.1) often gets ``trapped'' in
noisy areas or takes expensive detours, failing to find the true object boundary. Your manager asks you
to rewrite the core engine. The new requirement is:
\begin{itemize}[label=$\triangleright$]
\item ``The tool must guarantee finding the mathematically optimal boundary---the path with the absolute
minimum Total Cumulative Cost from Start to Target---even if it takes a bit more computation time.''
\end{itemize}
Please \textbf{trace the path} generated by this method from \textbf{S} to \textbf{T} and
\textbf{calculate the total cost.}
\begin{itemize}
\item Path Sequence: \underline{\hspace{4.5cm}} (e.g., (0,0) $\to$ (0,1) $\to$ \dots)
\item Total Path Cost: \underline{\hspace{4.5cm}}
\end{itemize}
\item {[2\%]} The solution in Q 4.2 works great on small grids. However, when applied to a 4K Image (with
\textbf{N} pixels), the tool becomes laggy. Profiling reveals the bottleneck is in the \textbf{Open Set
(Priority Queue)}---specifically, the operation of repeatedly finding and removing the pixel with the
minimum cumulative cost. You are \textbf{comparing two underlying data structures} for this Priority
Queue: a simple \textbf{Unsorted Array} vs. a \textbf{Binary Min-Heap}. Please analyze the Time
Complexity (Big-O) for the ``Extract-Min'' operation in both cases:
\begin{itemize}
\item \textbf{Case 1:} You use a simple \textbf{Unsorted Array} to store the candidate pixels, the time
complexity for Extract-Min is: \underline{\hspace{4.5cm}}
\item \textbf{Case 2:} You switch to a \textbf{Binary Min-Heap} to store the candidate pixels, the time
complexity for Extract-Min is: \underline{\hspace{4.5cm}}
\end{itemize}
\end{enumerate}
\end{pastexam}

\begin{pastsoln}
答案：(1) $(0,0) \to (0,1) \to (0,2) \to (1,2) \to (2,2) \to (2,3) \to (2,4) \to (3,4)$，成本 $27$；
(2) $(0,0) \to (1,0) \to (2,0) \to (2,1) \to (2,2) \to (2,3) \to (2,4) \to (3,4)$，成本 $22$；
(3) 未排序陣列 $O(N)$，二元最小堆積 $O(\log N)$。

\textbf{(1) 貪心走法。} 原題的「to get closer to the target」有兩種讀法：只考慮往右、往下（會更接近 T
的方向），或考慮所有沒走過的鄰居（不回頭）。 在這張圖上兩種讀法走出同一條路，下面列出所有沒走過的
鄰居（成本相同時照右、下、上、左的順序）：
\begin{center}
\small
\begin{tabular}{lll}
\toprule
目前位置 & 候選的鄰居（成本） & 走到 \\
\midrule
$(0,0)$ S & 右 $(0,1)$ 4、下 $(1,0)$ 7 & $(0,1)$ \\
$(0,1)$ & 右 $(0,2)$ 5、下 $(1,1)$ 9 & $(0,2)$ \\
$(0,2)$ & 右 $(0,3)$ 9、下 $(1,2)$ 8 & $(1,2)$ \\
$(1,2)$ & 右 $(1,3)$ 6、下 $(2,2)$ 4、左 $(1,1)$ 9 & $(2,2)$ \\
$(2,2)$ & 右 $(2,3)$ 1、下 $(3,2)$ 6、左 $(2,1)$ 2 & $(2,3)$ \\
$(2,3)$ & 右 $(2,4)$ 5、下 $(3,3)$ 7、上 $(1,3)$ 6 & $(2,4)$ \\
$(2,4)$ & 下 $(3,4)$ T 0、上 $(1,4)$ 9 & $(3,4)$ T \\
\bottomrule
\end{tabular}
\end{center}
成本 $4 + 5 + 8 + 4 + 1 + 5 + 0 = 27$。 第一步因為 $4 < 7$ 往右走進了上面那排，之後被 $9$、$8$ 這些高成本
的像素擋住，只好從 $(0,2)$ 往下穿過 $8$。

\textbf{(2) 最小成本路徑：Dijkstra。} 把每個像素當頂點，相鄰像素之間連邊；走進像素 $v$ 的成本是 $v$
的能量，所以邊 $u \to v$ 的權重是 $\mathit{cost}(v)$，都非負，可以用 Dijkstra（資結書第~25.2~節）。 從
S 出發，每次取出目前累積成本最小、還沒定案的像素，定案後鬆弛它的鄰居。 取出的順序（括號裡是定案的
距離）前段是
\begin{gather*}
(0,0)\,0,\ (0,1)\,4,\ (1,0)\,7,\ (0,2)\,9,\ (2,0)\,10,\ (2,1)\,12,\\
(1,1)\,13,\ (2,2)\,16,\ (3,0)\,16,\ (1,2)\,17,\ (2,3)\,17,\ \ldots
\end{gather*}
最後每個像素的最短累積成本：
\begin{center}
\begin{tabular}{c|ccccc}
 & 0 & 1 & 2 & 3 & 4 \\ \hline
0 & \textbf{0} & 4 & 9 & 18 & 26 \\
1 & \textbf{7} & 13 & 17 & 23 & 31 \\
2 & \textbf{10} & \textbf{12} & \textbf{16} & \textbf{17} & \textbf{22} \\
3 & 16 & 17 & 22 & 24 & \textbf{22} \\
\end{tabular}
\end{center}
T $= (3,4)$ 的距離是 $22$。 沿前驅往回走（粗體）：$(3,4) \leftarrow (2,4) \leftarrow (2,3) \leftarrow
(2,2) \leftarrow (2,1) \leftarrow (2,0) \leftarrow (1,0) \leftarrow (0,0)$，成本
$7 + 3 + 2 + 4 + 1 + 5 + 0 = 22$。 這條路先往下走進第 2 列那一排低成本的像素（$3, 2, 4, 1, 5$），比 (1)
省了 $5$。 例如 $(2,2)$ 的距離 $16$ 來自 $(2,1)$ 的 $12 + 4$，而貪心從上方 $(1,2)$ 走到 $(2,2)$ 時已經
累積了 $17 + 4 = 21$。 從 $(2,3)$ 往下 $(3,3) \to (3,4)$ 要 $7 + 0$，比往右 $(2,4) \to (3,4)$ 的 $5 + 0$
貴，所以最後一段走右邊。 最佳路徑只有這一條。

\textbf{(3)} 優先佇列裡最多有 $O(N)$ 個像素。
\begin{itemize}
\item 未排序陣列：插入直接放最後面 $O(1)$，但 Extract-Min 要掃過全部元素找最小值，再和最後一個交換後
刪掉，$O(N)$。
\item 二元最小堆積：最小值在根，取出後把最後一個元素放到根再下沉，$O(\log N)$（資結書第~8.3~節）。
\end{itemize}
每個像素最多四條邊，$E = O(N)$。 Dijkstra 要做 $N$ 次 Extract-Min：用陣列總共 $O(N^2)$；用堆積加上
$O(E)$ 次 decrease-key（每次 $O(\log N)$），總共 $O((N + E)\log N) = O(N \log N)$。 4K 影像
$N \approx 8.3 \times 10^6$，$N^2$ 和 $N \log N$ 差了幾十萬倍，這就是改用堆積的理由。
\end{pastsoln}
